\section{The Subscription That Already Worked}
\label{sec:ch14-the-subscription-that-already-worked}
\index{subscriptions!through a router|(}%

Before changing anything, check out tag \code{ch13}, start the seven services
and the router, and post a subscription to the router rather than to Reviews:

\begin{minted}[fontsize=\footnotesize,breaklines]{graphql}
subscription { onReviewAdded(productId: "UHJvZHVjdDoAAACgAAAAQIAAAAAAAAAW") {
  id rating body author { displayName } } }
\end{minted}

The router answers \code{200} with a content type of
\code{text/event-stream}, sends a heartbeat, and then goes quiet, which is what
a subscription with nothing to say looks like. Submit a review and a frame
arrives:

\begin{minted}[fontsize=\scriptsize,breaklines]{text}
event: next
data: {"data":{"onReviewAdded":{"id":"UmV2aWV3OjHunwGkClh8nA5sKCe2ci8=",
       "rating":1,"body":"Event 1.",
       "author":{"displayName":"Mateusz Kowalczyk"}}}}
\end{minted}

\index{subscriptions!a field the subgraph cannot fully answer}%
Read the last line twice. \code{author.displayName} is a field of
\code{Customer}, and \code{Customer} in the Reviews subgraph has been a
two-line stub since chapter~\ref{ch:strangling-the-monolith}: an identifier and
nothing else. Reviews does not know that customer's name and never did. So the
router did not simply forward what the subgraph pushed at it. It took the
subgraph's message apart, went to Accounts for the missing field, and put the
answer back together before the client saw a byte.

That is the shape of the whole chapter and it is worth being clear that nobody
built it. This is tag \code{ch13}, untouched.

\subsection*{What the composed config had been carrying}

\index{supergraph.json@\code{supergraph.json}!the subscription block}%
Chapter~\ref{ch:composition} pulled a router execution config apart field by
field and found a block it could not ask a question of. Every one of the seven
datasources carries this, with its own url:

\begin{minted}[fontsize=\footnotesize,breaklines]{json}
"subscription": {
  "enabled": true,
  "url": {"staticVariableContent": "http://localhost:5105/graphql"},
  "protocol": "GRAPHQL_SUBSCRIPTION_PROTOCOL_WS",
  "websocketSubprotocol": "GRAPHQL_WEBSOCKET_SUBPROTOCOL_AUTO"
}
\end{minted}

Including the six subgraphs that have no \code{Subscription} type at all.
Catalog will never be asked to hold a connection open and the composer wrote it
a subscription block anyway, because the block describes how the router would
reach that subgraph rather than whether it will ever need to.
Section~\ref{sec:ch14-two-protocols-one-name} is about the last line of it and
section~\ref{sec:ch14-five-transports} about the two before.

\subsection*{A plan with a node kind we have not seen}

\index{query plan!the \code{Trigger} node}%
Ask the router for the plan the way chapter~\ref{ch:enter-the-router} does,
with \code{X-WG-Include-Query-Plan} and \code{X-WG-Skip-Loader}, and the shape
is new:

\begin{minted}[fontsize=\scriptsize,breaklines]{text}
Sequence
  Trigger  reviews   fetchId 0
           subscription($a: ID!){ onReviewAdded(productId: $a){ id rating body
                                    author { __typename id } } }
  Single
    Entity accounts  fetchId 1  dependsOnFetchIds [0]
           query($representations: [_Any!]!){ _entities(representations: $representations){
                   ... on Customer { __typename displayName } } }
\end{minted}

Every plan in chapters~\ref{ch:how-a-federated-query-actually-runs}
and~\ref{ch:enter-the-router} through~\ref{ch:hard-modeling-problems} was built
from \code{Sequence}, \code{Parallel}, \code{Single} and \code{Entity}.
\code{Trigger} is a fifth kind and it exists because a subscription is the one
operation whose first step does not finish. The children hang off it in the
ordinary way, with the ordinary \code{dependsOnFetchIds}, and the ordinary
representation built from a key.

\index{query plan!what the trigger asks for}%
Notice what the trigger asks the subgraph for: \code{author \{ \_\_typename id
\}}. Not the display name. The planner knows Reviews cannot answer that, so it
does not ask, exactly as it would not ask in a query. What the subscription
field returns is a \code{Review}, and a \code{Review} is an entity that four
services contribute to; being pushed rather than pulled changes nothing about
how it gets assembled.

\subsection*{One write, and where it had never been sent}

\index{mutations!through the router}%
You cannot watch any of this without writing something, and
\code{submitReview} has been exercised against the Reviews subgraph directly
since chapter~\ref{ch:schema-design-that-survives-change}.
Chapter~\ref{ch:enter-the-router} recorded that as unfinished business and
chapter~\ref{ch:strangling-the-monolith} recorded it again. Sending it to the
router now:

\begin{minted}[fontsize=\scriptsize,breaklines]{text}
Sequence
  Single  reviews   mutation($a: SubmitReviewInput!){ submitReview(input: $a){
                      review { id rating author { __typename id } } } }
  Entity  accounts  dependsOnFetchIds [0]
                    ... on Customer { __typename displayName email }
\end{minted}

Two things fall out of that. The first is that a mutation's payload is resolved
across a boundary exactly like a query's, so \code{review.author.displayName}
comes back filled in, where the same mutation sent to Reviews directly answers
null for it. That null has been in the tests since
chapter~\ref{ch:strangling-the-monolith} made \code{Customer} a stub in this
service, and it was never a fact about the mutation. It was a fact about where
the request was sent.

The second is a difference between the two plans that I noticed and did not
chase: the mutation's plan carries a \code{normalizedQuery} member and the
subscription's does not. The research note records it as unmeasured with
chapter~\ref{ch:inside-the-router-and-the-composer} as its owner, because the
plan cache is that chapter's and four chapters have now asked it a different
question.

\index{errors!a typed error through the router}%
The typed errors survive the trip. Submit a review a customer has already
written and the payload comes back the way
chapter~\ref{ch:schema-design-that-survives-change} designed it:

\begin{minted}[fontsize=\footnotesize,breaklines]{json}
{"data":{"submitReview":{"review":null,"errors":[
  {"__typename":"DuplicateReviewError","message":"You have already reviewed this product."}]}}}
\end{minted}

One thing I got wrong on the way there is worth a sentence, because it looked
like a router defect for about ten minutes. Selecting \code{errors \{ message
\}} on that payload produces this:

\begin{minted}[fontsize=\footnotesize,breaklines]{text}
cannot select field: message on union: SubmitReviewError
\end{minted}

which reads like a planner limitation and is not one. \code{SubmitReviewError}
is a union, selecting a field on a union without a fragment is invalid GraphQL,
and the router is right. The correct selection is \code{errors \{ \_\_typename
... on Error \{ message \} \}}. A tool that disagrees with you is worth checking
against the specification before it is worth writing up.

\medskip
So the feature works and the write path works, and neither needed anything
built. The next question is what a message costs, which is where the interest
actually is.

\index{subscriptions!through a router|)}%
